\exercisesection{Exercises for § 4}

If \(\mathbb{R}\) is a root system, denote by \(W^{+}(R)\) the set of elements of \(W(R)\) of determinant 1.

¶ 1) Let R be a root system of type E8.

\begin{enumerate}
\def\labelenumi{\alph{enumi})}
\item
  Show that if \(\alpha, \beta \in \mathbb{R}\) are congruent modulo \(2\mathrm{Q}(\mathbb{R})\) , then \(\beta = \pm \alpha\) .
\item
  Deduce from a) that, if \(w \in \mathrm{W}(\mathrm{R})\) acts trivially on \(\mathrm{Q}(\mathrm{R}) / 2\mathrm{Q}(\mathrm{R})\) , then \(w = \pm 1\) .
\item
  With the notation of no. 10, show that the quadratic form \(\frac{1}{2}(x|x)\) on Q(R) defines by passage to the quotient a non-degenerate quadratic form \(q_{8}\) on the \(\mathbf{F}_2\) -vector space Q(R)/2Q(R). Show that the pseudo-discriminant of \(q_{8}\) (cf.~Algebra, Chap. IX, §9. Exerc. 9) is zero, and that \(q_{8}\) is of index 4.
\item
  Let \(\mathbf{O}(q_{8})\) be the orthogonal group of \(\mathrm{Q}(\mathrm{R})/2\mathrm{Q}(\mathrm{R})\) for this form. Define, by passing to the quotient, a homomorphism
\end{enumerate}

\[
h: \mathrm{W} (\mathrm{R}) \rightarrow \mathbf {O} (q _ {8}).
\]

Show (by comparing orders) that the sequence

\[
1 \rightarrow \{1, - 1 \} \rightarrow \mathrm{W(R)} \stackrel {h} {\longrightarrow} \mathbf {O} (q _ {8}) \rightarrow 1
\]

is exact.

\begin{enumerate}
\def\labelenumi{\alph{enumi})}
\setcounter{enumi}{4}
\tightlist
\item
  Show that the image of \(\mathrm{W}^{+}(\mathrm{R})\) under \(h\) is the subgroup \(\mathbf{O}^{+}(q_{8})\) of \(\mathbf{O}(q_8)\) defined in Algebra. Chap. IX, §9, no. 5. Deduce that \(\mathrm{W}^{+}(\mathrm{R}) / \{1, -1\}\) is a simple non-abelian group.
\end{enumerate}

\begin{enumerate}
\def\labelenumi{\arabic{enumi})}
\setcounter{enumi}{1}
\tightlist
\item
  Let R be a root system of type \(E_{6}\) .
\end{enumerate}

\begin{enumerate}
\def\labelenumi{\alph{enumi})}
\item
  Put \(E = Q(R)/3P(R)\) . This is an \(F_{3}\) -vector space of dimension 5. With the notation of no. 12, show that the scalar product \((x|y)\) defines on E a non-degenerate symmetric bilinear form \(\varphi\) . Show that any two distinct elements of R have distinct images in E.
\item
  Let \(\mathbf{O}(\varphi)\) be the orthogonal group of \(\varphi\) . Then, \(\mathbf{O}(\varphi) = \{1, -1\} \times \mathbf{SO}(\varphi)\) . The spinor norm defines a surjective homomorphism from \(\mathbf{SO}(\varphi)\) to \(\{1, -1\}\) whose kernel is denoted by \(\mathbf{SO}^{+}(\varphi)\) . The group \(\mathbf{SO}^{+}(\varphi)\) is simple of order 25920. The quotient \(\mathbf{O}(\varphi)/\mathbf{SO}^{+}(\varphi)\) is of type (2, 2). Deduce that \(\mathbf{O}(\varphi)\) contains a unique subgroup \(\Omega(\varphi)\) of index 2, distinct from \(\mathbf{SO}(\varphi)\) and not containing \(-1\) .
\item
  Every element of A(R) defines by passage to the quotient an element of \(\mathbf{O}(\varphi)\) . Show (by comparing orders) that this gives a homomorphism from A(R) to \(\mathbf{O}(\varphi)\) . The image of W(R) under this homomorphism is \(\Omega(\varphi)\) , and that of \(\mathrm{W}^{+}(\mathbb{R})\) is \(\mathbf{SO}^{+}(\varphi)\) . Hence, \(\mathrm{W}(\mathbb{R})\) is an extension of \(\mathbf{Z}/2\mathbf{Z}\) by a simple group of order 25920.
\item
  Let \(F = Q(R)/2Q(R)\) . This is an \(F_{2}\) -vector space of dimension 6. Show that the quadratic form \(\frac{1}{2}(x|x)\) defines by passage to the quotient a non-degenerate quadratic form \(q_{6}\) on F. of pseudo-discriminant equal to 1. If \(\mathbf{O}(q_{6})\) denotes the corresponding orthogonal group, define, by passing to the quotient, a homomorphism
\end{enumerate}

\[
h: \mathrm{W} (\mathrm{R}) \rightarrow \mathbf {O} (q _ {6}).
\]

Show that \(h\) is injective (note that \(-1 \notin \mathrm{W}(\mathbb{R})\) ), and then that it is an isomorphism (compare orders). Deduce that there is an isomorphism from \(\mathrm{W}^{+}(\mathbb{R})\) to \(\mathbf{O}^{+}(q_{6})\) (cf.~Algebra, Chap. IX, §9, no. 5).

\begin{enumerate}
\def\labelenumi{\alph{enumi})}
\setcounter{enumi}{2}
\tightlist
\item
  By comparing c) and d), show \(^{4}\) that \(\mathbf{SO}^{+}(\varphi)\) is isomorphic to \(\mathbf{O}^{+}(q_{6})\) .
\end{enumerate}

¶ 3) Let R be a root system of type \(E_{7}\).

\begin{enumerate}
\def\labelenumi{\alph{enumi})}
\item
  Put \(E = Q(R)/2P(R)\) . This is an \(F_{2}\) -vector space of dimension 6. With the notation in no. 11, show that the scalar product \((x|y)\) defines a non-degenerate alternating bilinear form on E.
\item
  Deduce from a) the existence of an exact sequence
\end{enumerate}

\[
1 \rightarrow \{1, - 1 \} \rightarrow \mathrm{W} (\mathrm{R}) \xrightarrow {h} \mathbf {S p} (6, \mathbf {F} _ {2}) \rightarrow 1.
\]

(Use the fact that \(\mathbf{Sp}(6,\mathbf{F}_2)\) is of order \(2^{9}.3^{4}.5.7\) .)

\begin{enumerate}
\def\labelenumi{\alph{enumi})}
\setcounter{enumi}{2}
\tightlist
\item
  Show that the restriction of \(h\) to \(\mathrm{W}^{+}(\mathrm{R})\) is an isomorphism from \(\mathrm{W}^{+}(\mathrm{R})\) to \(\mathbf{Sp}(6, \mathbf{F}_2)\) .
\end{enumerate}

\(d\) ) Give a second proof of \(b\) ) by using the quadratic form \(q_{7}\) on the \(\mathbf{F}_2\) -vector space \(\mathrm{Q(R) / 2Q(R)}\) induced by \(\frac{1}{2} (x|x)\) , as well as the isomorphism

\[
\mathbf {O} (q _ {7}) \rightarrow \mathbf {S p} (6, \mathbf {F} _ {2}).
\]

\begin{enumerate}
\def\labelenumi{\arabic{enumi})}
\setcounter{enumi}{3}
\tightlist
\item
  Let R be an irreducible reduced root system in V. \((\alpha_{1},\ldots,\alpha_{l})\) a basis of R. \(\tilde{\alpha}\) the highest root. Put \(\tilde{\alpha}=n_{1}\alpha_{1}+\cdots+n_{l}\alpha_{l}\) . We are going to determine the maximal closed, symmetric subsets of R, distinct from R.
\end{enumerate}

\begin{enumerate}
\def\labelenumi{\alph{enumi})}
\item
  Let \(i \in \{1, 2, \ldots, l\}\) . Let \(R_i\) be the set of \(\alpha \in R\) which are linear combinations of the \(\alpha_j\) for \(j \neq i\) . Show that \(R_i\) is maximal if and only if \(n_i = 1\) .
\item
  Let \(i \in \{1, 2, \ldots, l\}\) and assume that \(n_i > 1\) . Let \(S_i\) be the set of roots \(\sum_{j=1}^{l} m_j \alpha_j\) with \(m_i \equiv 0 \pmod{n_i}\) . Show that \(S_i\) is maximal if and only if \(n_i\) is prime. (If \(n_i = ab\) with a \textgreater{} 1, b \textgreater{} 1, consider the subset \(S'\) of R consisting of the roots \(\sum_{j} m_j \alpha_j\) with \(m_i \equiv 0 \pmod{a}\) , and show that \(S'\) strictly contains \(S_i\) .) Show that the roots \(-\tilde{\alpha}, \alpha_j (j \neq i)\) form a basis of \(S_i\) (which is of rank l). Deduce the Dynkin graph of \(S_i\) .
\item
  Every maximal closed, symmetric subset of R is transformed by an element of W(R) into one of the subsets described in a) or b). (Let \(\Sigma\) be such a subset. Then \(\Sigma = R \cap H\) , where H is a subgroup of Q(R) of finite index (§1, Exerc. 6 b)); we can assume that Q(R)/H is cyclic. Then there exists \(u^{*} \in V^{*}\) such that \(\Sigma\) is the set of \(\alpha \in R\) such that \(\langle u^{*}, \alpha \rangle \in Z\) .
\item
  List the maximal, closed, symmetric subsets of R for the different types of irreducible reduced root systems.
\end{enumerate}

\begin{enumerate}
\def\labelenumi{\arabic{enumi})}
\setcounter{enumi}{4}
\tightlist
\item
  Let R be a root system, and \(P'(R)\) the subgroup of \(P(R)\) introduced in §1. Exerc. 8. Show that \(P'(R) = P(R)\) if R is of type \(A_l\) with l even, or \(B_l\) with \(l \equiv 0,3 \pmod{4}\) , or \(D_l\) with \(l \equiv 0,1 \pmod{4}\) , or \(G_2\) , or \(F_4\) , or \(E_6\) , or \(E_8\) . Show that \(P'(R) = Q(R)\) if R is of type \(C_l\) , or \(B_l\) with \(l \equiv 1,2 \pmod{4}\) , or \(E_7\) , or \(A_l\) . If R is of type \(A_l\) with l odd and \textgreater1, \(P'(R)/Q(R)\) is the unique subgroup of index 2 in the cyclic group \(P(R)/Q(R)\) . If R is of type \(D_l\) with \(l \equiv 2,3 \pmod{4}\) , \(P'(R)/Q(R)\) is the unique subgroup of order 2 of \(P(R)/Q(R)\) stable under \(\Lambda(R)\) .
\end{enumerate}

¶ 6) Let R be an irreducible reduced root system, \((\alpha_{1},\ldots,\alpha_{l})\) a basis of R, \(p_{1}\alpha_{1}+\cdots+p_{l}\alpha_{l}\) the highest root, \(a_{1}\alpha_{1}+\cdots+a_{l}\alpha_{l}\) the sum of the positive roots, \(m_{1},\ldots,m_{l}\) the exponents of W(R), g its order, f the connection index.

\begin{enumerate}
\def\labelenumi{\alph{enumi})}
\tightlist
\item
  Verify that, in each case,
\end{enumerate}

\[
l! p _ {1} p _ {2} \dots p _ {l} m _ {1} m _ {2} \dots m _ {l} = a _ {1} a _ {2} \dots a _ {l}.
\]

\begin{enumerate}
\def\labelenumi{\alph{enumi})}
\setcounter{enumi}{1}
\tightlist
\item
  Show that
\end{enumerate}

\[
g m _ {1} m _ {2} \dots m _ {l} = f a _ {1} a _ {2} \dots a _ {l}.
\]

(Use \(a\) ) and Prop. 7 of §2, no. 4.)

\begin{enumerate}
\def\labelenumi{\alph{enumi})}
\setcounter{enumi}{2}
\tightlist
\item
  For any positive root \(\alpha = \sum_{i} c_{i} q_{i}\) , let \(c(\alpha) = \sum_{i} c_{i}\) . Calculate in each case the polynomial
\end{enumerate}

\[
\mathrm{P} (t) = \sum_ {\alpha > 0} t ^ {c (\alpha)}.
\]

Verify \(^{5}\) that \(P(t) = t \sum_{i=1}^{l} (1 + t + \cdots + t^{m_{i}-1})\) .

\begin{enumerate}
\def\labelenumi{\arabic{enumi})}
\setcounter{enumi}{6}
\tightlist
\item
  Let R be an irreducible reduced root system.
\end{enumerate}

\begin{enumerate}
\def\labelenumi{\alph{enumi})}
\item
  Verify that the canonical homomorphism from \(A(R)/W(R)\) to the group of automorphisms of \(P(R)/Q(R)\) is injective.
\item
  Deduce that -1 belongs to W(R) if and only if Q(R) ⊆ 2P(R).
\end{enumerate}

\begin{enumerate}
\def\labelenumi{\arabic{enumi})}
\setcounter{enumi}{7}
\item
  Let \(R_{1}\) and \(R_{2}\) be two irreducible reduced root systems. Show that if \(W(R_{1})\) and \(W(R_{2})\) have the same order. \(R_{1}\) is isomorphic to \(R_{2}\) or \(R_{2}^{-}\) . (Use the classification.) Does this result still hold without the assumption that \(R_{1}\) is irreducible?
\item
  Let R be a root system of rank l, and let p be a prime number dividing the order of A(R). Show that \(p \leqslant l + 1\) . (Reduce to the irreducible case, and use the classification.)
\item
  \begin{enumerate}
  \def\labelenumii{\alph{enumii})}
  \tightlist
  \item
    Let \((W, S)\) be an irreducible, finite Coxeter system. Put
  \end{enumerate}
\end{enumerate}

\[
W (t) = \sum_ {w \in W} t ^ {l (w)} \quad (\text { cf.   Chap.   IV,   } \S 1, \text {   Exerc.   26. })
\]

Let \(m_{1},\ldots ,m_{l}\) be the exponents of W (Chap. V, §6. no. 2). Verify the formula

\[
\mathrm{W} (t) = \prod_ {i = 1} ^ {l} \left(1 + t + \dots + t ^ {m _ {i}}\right)
\]

for small values of \(l\) (use Th. 1 and Exerc. 26 of Chap. IV, § 1) \(^{6}\) .

\begin{enumerate}
\def\labelenumi{\alph{enumi})}
\setcounter{enumi}{1}
\tightlist
\item
  Let R be a reduced, irreducible root system and \(W\) (resp. \(W_{a}\)) the Weyl group (resp. the affine Weyl group) of R, equipped with the Coxeter group structure determined by the choice of an alcove. Define \(W(t)\) and the exponents \(m_{i}\) as above and put
\end{enumerate}

\[
W _ {a} (t) = \sum_ {w \in W _ {a}} t ^ {l (w)}.
\]

Verify the formula

\[
\mathrm{W} _ {a} (t) = \mathrm{W} (t) \prod_ {i = 1} ^ {l} \frac {1}{1 - t ^ {m _ {i}}} = \prod_ {i = 1} ^ {l} \frac {1 + t + \cdots + t ^ {m _ {i}}}{1 - t ^ {m _ {i}}}
\]

for small values of \(l\) (use Th. 2 and Exerc. 26 of Chap. IV, § 1) \(^{7}\) .

¶ 11) Let (W.S) be a Coxeter system of type \(H_{3}\) (cf.~Th. 1).

\begin{enumerate}
\def\labelenumi{\alph{enumi})}
\item
  With the notation of Chap. V, §6, no. 2. Proof of Lemma 2, show that \((z', z'') = \pi/5\) ; deduce that the Coxeter number \(h\) of \(W\) is equal to 10, and that the exponents of \(W\) are 1, 5 and 9.
\item
  Show, by using a), that \(\operatorname{Card}(W) = 120\) , and that the number of reflections in W is 15.
\item
  Recover the formula \(\operatorname{Card}(W) = 120\) by applying Exerc. 5 of Chap. V. §3. d) Let \(\mathcal{A}_5\) be the alternating group of \(\{1, \ldots, 5\}\) ; if \(a, b, c, d\) are distinct elements of \(\{1, \ldots, 5\}\) , denote by \((ab)\) the transposition of \(a\) and \(b\) , and \((ab)(cd)\) the product of the transpositions \((ab)\) and \((cd)\) . Let
\end{enumerate}

\[
r _ {1} = (14) (23), \quad r _ {2} = (12) (45), \quad r _ {3} = (12) (34).
\]

Show that \((r_1r_2)^5 = (r_2r_3)^3 = (r_1r_3)^2 = 1\) . Deduce the existence of a homomorphism \(f: \mathrm{W} \to \mathcal{A}_5\) that takes S to \(\{r_1, r_2, r_3\}\) ; show that \(f\) is surjective.\\
e) Let \(\varepsilon: \mathrm{W} \to \{\pm 1\}\) be the homomorphism \(w \mapsto (-1)^{l(w)}\) . Show that

\[
(f, \varepsilon): W \rightarrow \mathcal {A} _ {5} \times \{\pm 1 \}
\]

is an isomorphism. (Use the fact that the two groups being considered have the same order.)

¶ 12) Let \((1, i, j, k)\) be the canonical basis of the field H of quaternions, by means of which H is identified with \(\mathbf{R}^4\) . Equip H with the scalar product \(\frac{1}{2}(x\bar{y} + y\bar{x})\) . Let \(\Gamma\) be the multiplicative group of quaternions of norm 1.

\begin{enumerate}
\def\labelenumi{\alph{enumi})}
\item
  If \(a \in \Gamma\) , the orthogonal reflection \(s_{a}\) in H which transforms a to -a is the map \(x \mapsto -a\bar{x}a\) .
\item
  Let \(q = \cos \frac{\pi}{5} - \frac{1}{2}i + (\cos \frac{3\pi}{5})j \in \Gamma\) , and \(r = \frac{1}{2}(1 + i + j + k) \in \Gamma\) . Let Q be the set of quaternions obtained from 1, q, r by arbitrary even permutations and sign changes of the coordinates. Then Q is a subgroup of \(\Gamma\) of order 120.
\item
  Let W be the subgroup of \(\mathbf{GL}(\mathbf{H})\) generated by the \(s_{a}\) for \(a \in Q\) . Show that W leaves Q stable, is finite, irreducible, and non-crystallographic: deduce that W is of type \(H_{4}\) (cf.~Theorem 1).
\item
  Show that W acts transitively on Q (use Prop. 3 of Chap. IV, §1).
\item
  Let \(a_0 \in Q\) , let \(V\) be the vector subspace of \(\mathbf{H}\) orthogonal to \(a_0\) , and let \(W_0\) be the stabiliser of \(a_0\) in \(W\) . Show that the restriction of \(W_0\) to \(V\) is an irreducible group generated by reflections (Chap. V. §3, no. 3) and that it is non-crystallographic. Deduce that \(W_0\) is a Coxeter group of type \(H_3\) .
\item
  Show that \(\operatorname{Card}(W)=2^{6}3^{2}5^{2}\) (use (d), (e) and Exerc. 11).
\end{enumerate}

g) Show that the Coxeter number of \(\mathbf{W}\) is equal to 30, and that its exponents are 1, 11, 19, 29.

\begin{enumerate}
\def\labelenumi{\arabic{enumi})}
\setcounter{enumi}{12}
\tightlist
\item
  Let V be the hyperplane in \(R^{9}\) with equation \(x_{1} + \cdots + x_{9} = 0\) . Let R be the subset of V consisting of the points
\end{enumerate}

\[
(2. 2, 2. - 1, \dots 1. - 1, \dots 1, \dots 1, \dots 1), \quad (- 2. - 2. - 2. 1, 1. 1. 1. 1. 1),
\]

\[
(3. - 3. 0. 0. 0. 0. 0. 0. 0)
\]

together with the points obtained from these by permuting the coordinates. Show that R is a root system in V of type \(E_{8}\) .

\begin{enumerate}
\def\labelenumi{\arabic{enumi})}
\setcounter{enumi}{13}
\item
  With the notation in no. 7, show that the automorphism of \(\mathbf{R}^{l+1}\) which transforms \(\varepsilon_1\) to \(\varepsilon_2\) , \(\varepsilon_2\) to \(\varepsilon_3\) , \(\ldots\) , \(\varepsilon_l\) to \(\varepsilon_{l+1}\) and \(\varepsilon_{l+1}\) to \(\varepsilon_1\) induces a Coxeter transformation of the system R of type \(A_l\) .
\item
  Determine the minuscule weights (§ 1. Exerc. 24) for each type of reduced, irreducible root system. (One finds the fundamental weights \(\omega_{1},\ldots,\omega_{l}\) for \(A_{l}\) , the weight \(\omega_{l}\) for \(B_{l}\) , the weight \(\omega_{1}\) for \(C_{l}\) , the weights \(\omega_{1},\omega_{l-1},\omega_{l}\) for \(D_{l}\) , the weights \(\omega_{1}\) and \(\omega_{6}\) for \(E_{6}\) , the weight \(\omega_{7}\) for \(E_{7}\) , and none for \(E_{8},F_{4}\) and \(G_{2}\) .)
\end{enumerate}

¶ 16) Let (W.S) be an irreducible, finite Coxeter system, and let n = Card(S).

\begin{enumerate}
\def\labelenumi{\alph{enumi})}
\item
  If \((W, S)\) is not of type \(F_4\) , show that there exists a subset \(X\) of \(S\) with \(n - 1\) elements such that \((W_X, X)\) is of type \(A_{n-1}\) .
\item
  Identify W with a subgroup of \(\mathbf{GL}(\mathbf{R}^{S})\) by means of the canonical representation (Chap. V, §4). Show that there exists a basis \((c_{1},\ldots,c_{n})\) of \(R^{S}\) such that, for every permutation \(\sigma\in S_{n}\) , the automorphism of \(R^{S}\) that transforms \(e_{i}\) to \(c_{\sigma(i)}\) for \(1\leqslant i\leqslant n\) belongs to W (``Burnside's Theorem''). (When (W.S) is not of type \(F_{4}\) , use a): when it is of type \(F_{4}\) , remark that W contains a subgroup of type \(D_{4}\) (cf.~no. 9), which reduces the problem to the preceding case.)
\item
  Let E be a subgroup of the group of automorphisms of (W.S). Show that the semi-direct product E.W of E by W embeds in a canonical way in \(\mathbf{GL}(\mathbf{R}^{\mathrm{S}})\) . Show that the subgroup of \(\mathbf{GL}(\mathbf{R}^{\mathrm{S}})\) thus defined is generated by reflections, except in the following four cases:
\end{enumerate}

\begin{enumerate}
\def\labelenumi{(\roman{enumi})}
\item
  (W.S) is of type \(A_{n}\) . \(n \geqslant 4\) : the group \(E\) is of order 2.
\item
  (W.S) is of type \(\mathrm{D}_4\) ; the group \(E\) is of order 3.
\item
  (W.S) is of type \(F_4\) ; the group \(E\) is of order 2.
\item
  (W.S) is of type \(\mathbf{E}_6\) ; the group \(E\) is of order 2.
\end{enumerate}

Show that, in cases (i) and (iv). E.W = \(\{\pm1\} \times W\) . Show that, in case (iii), the group E.W does not leave stable any lattice in \(R^{S}\) .

\begin{enumerate}
\def\labelenumi{\alph{enumi})}
\setcounter{enumi}{3}
\tightlist
\item
  Let \(W_{1}\) be a finite subgroup of \(\mathbf{GL}_{n}(\mathbf{R})\) generated by reflections, and assume that \(W_{1}\) is irreducible and essential. Let G be a finite subgroup of \(\mathbf{GL}_{n}(\mathbf{R})\) containing \(W_{1}\) . Show that either G is generated by reflections or is of the form E.W. where E and W are one of the types (i), (ii), (iii), (iv) of c).
\end{enumerate}

(Let W be the group generated by the reflections belonging to G. The group G permutes the chambers relative to W. Deduce, as in §2, no. 3, that G is of the form E.W as above.)
% semantic-fragments: legacy-input-seam
\relax{}
